\begin{proposition}[A global identification certificate for pooled moments]\label{prop:poolcertificate}
Suppose $K$ is known and the normalized moments are $t_{d,m}=\sum_{i=1}^K W_{di}s_i^m$, with distinct ordered real nodes and strictly positive weights, observed for $d=1,\ldots,D$ and $m=0,\ldots,L$, $L\ge K$. Form the $D(L-K+1)\times K$ block-Hankel matrix and vector
\[
H_{(d,j),\ell}=t_{d,j+\ell},\qquad b_{(d,j)}=-t_{d,j+K},
\quad 0\le j\le L-K,\quad0\le\ell<K.
\]
If $\operatorname{rank}H=K$, the moments identify all nodes and weights globally, up to relabelling, among populations with $K$ distinct nodes. At $L=l^*=K-1+\lceil K/D\rceil$, this condition holds outside a proper algebraic subset of the parameter space. On any compact ordered class with separated nodes, bounded weights, and $\operatorname{rank}H=K$ throughout, the moment map has a Lipschitz inverse and injective differential. On a class with nonempty interior, the global observation-cost exponent in Theorem~\ref{thm:pool} is therefore $2l^*$.
\end{proposition}

\begin{proof}
Let $P(z)=\prod_i(z-s_i)=z^K+\sum_{\ell=0}^{K-1}c_\ell z^\ell$. Multiplying $P(s_i)=0$ by $W_{di}s_i^j$ and summing gives $Hc=b$. Full column rank makes $c$ unique. Any other $K$-node population with the same moments has a monic annihilating polynomial satisfying the same system, so it has the same polynomial and the same roots. The first $K$ moments in each market then determine its weights through the nonsingular Vandermonde matrix. This is a uniqueness argument, without an optimization search.

For genericity let $q=\lceil K/D\rceil$ and $L=K-1+q$. With $V_{i,\ell}=s_i^\ell$, one has $H=BV$, where $B_{(d,j),i}=W_{di}s_i^j$, $0\le j<q$. Assign each node to one of $D$ groups of size at most $q$, give it weight one in its assigned market and zero elsewhere, and take distinct nodes. The columns of $B$ then form disjoint Vandermonde blocks, each of full column rank. Thus $H$ has rank $K$. A sufficiently small positive perturbation of every zero weight preserves a nonzero minor and gives strictly positive weights. The relevant minors are polynomials, so rank failure is a proper algebraic condition. This proves genericity; it does not assert full rank at every positive configuration.

For stability of the inverse, $c=(H^\top H)^{-1}H^\top b$ is smooth wherever $H$ has full column rank. Roots of $P$ are locally smooth functions of its coefficients when they are simple, and each market's weights are smooth functions of the ordered roots and its first $K$ moments. These maps give a local smooth left inverse of the moment map, hence an injective differential. On the stated compact class, the smallest singular value of $H$ and the node gaps are bounded away from zero. Local inverse bounds, compactness and injectivity give a global Lipschitz inverse: otherwise pairs with a vanishing image-to-parameter distance ratio would either have distinct limits, contradicting injectivity, or a common limit, contradicting a local inverse bound. The upper rate follows by the rescaled-moment argument of Theorem~\ref{thm:pool}(ii); its lower rate holds on the class's interior by part (i).
\end{proof}

This certificate is the multi-channel annihilating-polynomial method underlying Prony reconstruction, not a new signal-processing principle. It supplies a direct condition for the financial pooling model and closes the generic global-uniqueness question at $l^*$. Rank-deficient exceptional classes, unknown numbers of types, misspecified shared memories, cross-impact, and a robust numerical rank decision under noise require further treatment. The smallest singular value matters: full rank in exact arithmetic alone is not a guarantee of useful finite-sample precision. The proof is on paper; exact rational tests exercise the recurrence and its reconstruction for $2\le K\le8$ and $1\le D\le K$.
